\section{Discussion and Conclusions}\label{sec:conclusions}

Self-consistency in sense~(a) is achievable and is a constraint worth accepting.
Every capture efficiency in \cref{eq:Zsk} and every sink strength in
\cref{eq:KL}, for all nine families, is generated from the same diffusion
tensors that transport the mobile species, through one anisotropy parameter
$p^{m}$ per species. The formulation therefore cannot simultaneously hold that
diffusion is isotropic and that its sinks are biased by the anisotropy of that
diffusion, which is a state most cluster-dynamics treatments of zirconium are
free to occupy. The price is that anisotropy can no longer absorb a shortfall by
refitting a bias --- and that price turns out to be worth paying, because it
lets $p^{m}$ be set from transport rather than from loop counts. With
$p^{v}=1.02$ and $p^{i}=0.70$ fixed as inputs, the basal diameter comes out at
$100.7\,$nm against $95\,$nm measured without being steered there, and it stays
within $85$--$131\,$nm at every stage of the search. The prismatic families
follow on the same fixed anisotropy: density $0.88\times$ the measurement over
71 points and diameter $1.21\times$ over 44, across twelve temperatures and four
decades of dose, with a scatter of $10^{0.49}=3.1$ that is comparable to the
spread among the measurements themselves. A model that cannot hide a deficiency
in a bias factor is one whose deficiencies stay visible, and one whose
agreements mean something when they occur --- the basal diameter was a factor
of two low before the anisotropy was fixed, and no lever was applied to it.

The second sense holds where it can be made exact and is measured where it
cannot. Stored defects transfer to machine precision and loop number to the
integer draw, at every dose tested. The sink strength --- which the transfer does
not force --- came within $2\%$ of continuity only after an inconsistency in the
\cloop{} Burgers vector between the two representations was removed, and locating
it is what the ledger of \cref{eq:ledger} is for: a discontinuity there is a
crystallography error announcing itself, not a numerical tolerance to be widened.
What sense~(b) does not yet deliver on this geometry is feasibility. A quarter of
the basal loops on a \SI{500}{\nano\meter} grain do not fit inside it
(\cref{sec:results-discrete}), so a hybrid calculation there would have to carry
the refused population in the continuum and account for it separately --- and
that same constraint is why the basal diameter was calibrated against the
measurement rather than against a transfer factor.

Beyond what the two consistency constraints buy, spatial resolution changes the
answer and not merely its presentation. The boundary layers of
\cref{sec:results-mobile} are $8$--$46\,$nm wide, they order themselves by
species mobility, and they narrow by a factor of two as the sink density grows,
none of which a mean-field $k^{2}_{gb}$ can express. More consequentially, the
loop population in the outermost shell is numerous and tiny and it dominates
every domain average: without the boundary loop channel the domain and interior
means of the \cloop{} density differ by a factor of $114$. A spatially averaged
model does not have this problem because it cannot see it, and it reports the
wrong number without any indication that it has.

That resolution is affordable, and the reason is worth separating from the
physics. $2.1\times10^{8}$ stiff ordinary differential equations were integrated
in about an hour on a workstation, not because a million-equation stiff system
was solved but because the frozen-mobile construction turns it into $72\,928$
independent $36$-dimensional ones --- $6.5\times10^{9}$ fewer flops per
factorization --- while the genuinely coupled part, having no time derivative, is
solved eight times rather than thousands. The bill arrives as a first-order
splitting error whose dial is the fast-solve cadence, and the two sides of the
calculation cost within $7\%$ of each other, which is the balance that cadence
was chosen to produce.

With the split in place, conservation can be audited channel by channel, and it
closes once the boundary is carried explicitly. Both atom balances close on
production to the digit, with the loop-at-boundary channel taking $5.1\%$ of
vacancy production against $0.08\%$ of interstitial. That the channel had to be
booked separately rather than left in the residual generalizes: closure by
difference is not a check, because a channel omitted from it reappears silently
in the remainder and nothing looks wrong. The channel that had to be added to
close that balance is also the one that makes a prediction. The size-gated
removal of \cref{eq:gbloop} produces a denuded zone $\sim5\,$nm wide in which the
loop density falls to the concentration floor and the mean diameter recovers
monotonically from $6.8$ to $144\,$nm, and it demonstrates something a mean-field
treatment structurally cannot: a channel acting only within a diffusion length of
a surface raises the \cloop{} density in the deep interior by a factor of 4--5,
because removing the shell's sinks changes the mobile field everywhere.

What the calculation does not settle is the basal density --- $2.1\times$ the
measurement --- and the shape of that residual is more useful than its size. The
model carries more \cloop{} loops than are measured while their diameter now
agrees, which is a weaker and more specific statement than the earlier
calibrations supported, where number and size were wrong together: number and size are
linked, the number balance pinning total content and the growth balance pinning
size, so a model removing loops faster than it grows them has survivors that are
too young, and every lever exercised here moves along that constraint rather than
across it. Raising coalescence brings the density down and collapses the diameter
with it; the small-size floor moves the prismatic families and not the basal one;
gating coalescence to the large end is worse still, the size bias being already
inside the Avrami kernel so that applying it twice shrinks the family, weakens
the rate and \emph{raises} the density. Fixing $p^{m}$ from transport sharpens
the diagnosis rather than removing it: what now rests on its bound is the basal
capture prefactor, $Z^{0}_{v}$ and $\xi_{c}$ together, which is a statement about
absorption and cannot be mistaken for one about diffusion. The mechanism the data
ask for removes vacancy loop \emph{number} without shortening the survivors'
lives --- unfaulting to a glissile configuration that escapes to the network, a
dissolution channel for the perfect basal state, or an absorption cross-section
growing faster with size than $R$ --- and identifying it is the single most
valuable next piece of physics. The prismatic families carry a smaller version of
the same question: their density is reproduced in the mean, $0.88\times$, and
low by a factor of six near $600\,$K, a temperature trend no lever in the
calibrated set corrects. \Cref{sec:atrend} sets out the two readings --- a
nucleation term with no temperature dependence of its own, or high-temperature
data that do not measure the interstitial population --- and the
character-resolved measurement at $600$--$700\,$K that would separate them.

Two further reservations attach to the calibration. The basal chain is
uncalibrated and switched off, because nothing in the data supplies the collapse
and unfaulting barriers of \cref{eq:basal-chain}; switched on with placeholder
rates it inverts the family ordering, which is a property of the placeholders and
not a prediction. Its terminal state $c_p$ also has no number-loss channel beyond
coalescence --- it cannot unfault further and peripheral emission acts on content
only --- so under a sustained vacancy supply it is bounded only by the transfer
feeding it, and whether it should have a sink of its own must be settled before
$\nu_{\mathrm{uf}}$ is fitted, the two not being separately identifiable from a
steady $c_p$ density. And the \cloop{} data are thin: five points at two
temperatures, all above 14\,dpa, while the model places resolvable basal loops
from $10^{-2}\,$dpa with nothing constraining that extrapolation. No posterior
was sampled either; \cref{sec:objective} establishes that the reported set is a
maximum \emph{a posteriori} point under a stated likelihood and bounded uniform
priors, without a credible interval to accompany it.

Four steps follow, in the order their results feed one another.

\begin{enumerate}[itemsep=3pt]
\item Refit with the boundary channel active, on the spatially resolved model
      rather than in $0$-D, so the calibration and the calculation describe the
      same physics. The transfer factors of \cref{sec:slack} are a stopgap for
      exactly this, and measuring them on the present march would replace them.
\item Constrain $p^{m}$ atomistically. Because it is the only anisotropy the
      formulation has, a computed value enters \cref{eq:Emsplit} directly and
      either confirms \cref{eq:pinned-pm} or replaces it, and the interstitial
      migration split of $0.106\,$eV it implies at $573\,$K is the specific
      quantity to compute.
\item Run the Bayesian pipeline. At 0.2--0.5\,s per objective evaluation, a
      space-filling design screened by Sobol' indices, a Gaussian-process
      emulator and Markov-chain sampling of \cref{eq:posterior} are all
      affordable, and would replace the point estimate of \cref{tab:fitted} with
      a posterior, identify which parameters the data determine, and quantify how
      far the \cloop{} predictions can be trusted on five measurements.
\item Close the loop to dislocation dynamics. The conversion of
      \cref{sec:conversion} is one-way: it exports a population and does not feed
      it back into a solve. Transferring a family when \cref{eq:dcoarsen} fires,
      removing it from the continuum and continuing with the loops as discrete
      objects whose stress fields and climb are solved is what would make the
      model hybrid rather than continuum with an export. The obstacle is
      geometric and is measured in \cref{sec:results-discrete}.
\end{enumerate}

\section*{Acknowledgments}

This work was performed under a Statement of Work with Fluor Marine Propulsion,
LLC. The authors thank colleagues at Knolls Atomic Power Laboratory and Bettis
Atomic Power Laboratory for discussions on the experimental microstructure of
irradiated zirconium alloys.
